\chapter{Conclusion}
\hspace{1em}The limits exposed by the D-gene positive control point to the experimental design itself, not only to the platform's methodology. Both case studies compare a single accession against a small group of others using a binary phenotype call -- juicy versus dry, or high versus low TAA -- so each candidate gene list rests on very few data points and cannot distinguish a true causal signal from natural variation between individual plants. A study with more accessions per phenotype class, and a phenotype measured on a continuous scale rather than as a binary split, would let the causal locus be mapped statistically (for example, via quantitative trait locus, or QTL, mapping) rather than inferred from a handful of accession contrasts, giving each omics layer a stronger and more precise signal to work with.

\hspace{1em}Concretely, this means measuring TAA accumulation as an actual concentration across many accessions rather than sorting a few accessions into "high" and "low" bins, so that variation in TAA level can be linked directly to genotype through QTL mapping. The same applies to stem juiciness: rather than a binary juicy/dry classification, the degree of aerenchyma formation and pith parenchyma cell death could be measured directly and quantitatively, since this is the developmental process the D-gene actually controls. Both changes would let future work replace the coarse, small-panel comparisons used in this thesis with a properly powered genetic mapping study, and are the most direct route to strengthening the platform's candidate gene lists for both traits.
